\begin{proof}[Proof of \cref{obj:thm:rectangular-upper}]
Fix the public parameters in the statement. Every constant constructed below depends only on these parameters.

\begin{enumerate}
\item Fix admissible \(n,m,h,J\) and \(P\in\mathcal M\). Define
\[
N=n+m,\qquad S=\alpha+\beta,\qquad
\delta=\frac hJ,\qquad
q=\binom{d+2}{2},\qquad k=qJ^d.
\]
In particular, \(N\ge n\ge2\) and \(\delta>0\).
By \cref{obj:lem:local-polynomial-projection-facts}, there is a constant
\(0<C_{\mathrm{loc}}<\infty\), uniform over these choices, and a coefficient vector
\(\vartheta\in\mathbb R^q\) such that
\[
p_\vartheta(x)=r(x)^\top\vartheta,\qquad
\sup_{x\in C_h}|\tau_P(x)-p_\vartheta(x)|
 \le C_{\mathrm{loc}}h^\gamma,\qquad
\|\vartheta\|_2\le C_{\mathrm{loc}},
\qquad
r_0^\top\vartheta=\tau_P(x_0).
\]
The same result supplies
\[
\|r_0\|_2\le C_{\mathrm{loc}},\qquad
\|e_Pr_u-\Pi_J(e_Pr_u)\|_{L_2(\nu_h)}
 \le C_{\mathrm{loc}}\delta^\alpha,\qquad
\|\mu_{0,P}-\Pi_J\mu_{0,P}\|_{L_2(\nu_h)}
 \le C_{\mathrm{loc}}\delta^\beta.
\]
Here \(x_0\in C_h\), so the bound on \(r_0\) follows from the stated uniform coarse-basis bound.

The finite Borel basis functions and the finite moment sums in
\cref{obj:def:moments} make both empirical moments finite Borel functions of the dataset.
The eigenvalue cutoff is Borel: its matrix superlevel set is the intersection of the closed unit-vector quadratic-form inequalities. Matrix inversion on the invertible branch and clipping are Borel. Thus the decision in
\cref{obj:def:estimator}, with value zero on the failed-cutoff branch, is finite and Borel on every dataset. Composing with the dataset coordinate gives a Borel decision on the original dataset and randomizer.%
% lean: CausalSmith.Stat.TwosamplePointcateAnnotationFrontier.tHat_measurable

\item Define the population moments and the localized Taylor remainder by
\[
\bar R=\mathbb E_{\mathsf E_{n,m}(P)}\widehat R_{h,J},
\qquad
Q_{P,h,J}=\mathbb E_{\mathsf E_{n,m}(P)}\widehat Q_{h,J},
\qquad
g_{\mathrm{rem}}:C_h\longrightarrow\mathbb R,\quad
g_{\mathrm{rem}}(x)=\tau_P(x)-p_\vartheta(x).
\]
All localized integrals below are over \(C_h\).

The observed-response identity and the exchangeability condition in
\cref{obj:def:primitive-class} give, almost surely,
\[
\mathbb E_P[AY\mid X]=e_P(X)\mu_{1,P}(X),
\qquad
\mathbb E_P[Y\mid X]
 =\mu_{0,P}(X)+e_P(X)\tau_P(X).
\]
Indeed, \(AY=AY(1)\), and conditional exchangeability factors its conditional mean. Applying the same argument to \((1-A)Y(0)\) and adding gives the second identity.

The role counts in \cref{obj:def:roles} satisfy
\[
\ell=\lfloor n/2\rfloor\ge1,\qquad
t=N-\ell\ge1.
\]
Their records are independent under the sampling law in
\cref{obj:ass:sample-law}. Consequently, taking expectations in
\cref{obj:def:moments}, using uniform design and factoring each rectangular product, yields
\[
\begin{aligned}
\bar R_u
&=\int r_u e_P\mu_{1,P}\,d\nu_h
 -\int e_Pr_u\,\Pi_J(\mu_{0,P}+e_P\tau_P)\,d\nu_h,\\
(Q_{P,h,J})_{uv}
&=\int r_u e_Pr_v\,d\nu_h
 -\int e_Pr_u\,\Pi_J(e_Pr_v)\,d\nu_h.
\end{aligned}
\]
The response in the rectangular correction is \(Y\), so its conditional mean is the second conditional mean displayed above.

For completeness, writing the orthonormal fine basis from
\cref{obj:def:projection} as
\[
\mathbf b_J=(b_{J,\iota})_{\iota=1}^k,
\]
its finite expansion gives, for arbitrary
\(f_1,f_2\in L_2(\nu_h)\),
\[
\int f_1\Pi_Jf_2\,d\nu_h
=
\sum_{\iota=1}^k
 \left(\int b_{J,\iota} f_1\,d\nu_h\right)
 \left(\int b_{J,\iota} f_2\,d\nu_h\right)
=
\int (\Pi_Jf_1)f_2\,d\nu_h.
\]
Orthonormality also gives
\[
\int f_1f_2\,d\nu_h-\int f_1\Pi_Jf_2\,d\nu_h
=
\int(f_1-\Pi_Jf_1)(f_2-\Pi_Jf_2)\,d\nu_h,
\]
and
\[
\|\Pi_Jf_1\|_{L_2(\nu_h)}^2
+\|f_1-\Pi_Jf_1\|_{L_2(\nu_h)}^2
=\|f_1\|_{L_2(\nu_h)}^2.
\]
These identities show that the population matrix displayed above is symmetric, so it agrees with the expectation of the symmetrized empirical matrix.

Multiplying its entries by \(\vartheta\), substituting
\(\mu_{1,P}=\mu_{0,P}+\tau_P\), and using the residual identity gives
\[
\begin{aligned}
\bar R_u-(Q_{P,h,J}\vartheta)_u
&=\int
 \bigl(e_Pr_u-\Pi_J(e_Pr_u)\bigr)
 \bigl(\mu_{0,P}-\Pi_J\mu_{0,P}\bigr)\,d\nu_h\\
&\quad+\int r_u e_Pg_{\mathrm{rem}}\,d\nu_h
 -\int \Pi_J(e_Pr_u)e_Pg_{\mathrm{rem}}\,d\nu_h.
\end{aligned}
\]
Cauchy--Schwarz and the two residual bounds from the preceding step imply
\[
\left|\int
 \bigl(e_Pr_u-\Pi_J(e_Pr_u)\bigr)
 \bigl(\mu_{0,P}-\Pi_J\mu_{0,P}\bigr)\,d\nu_h\right|
\le C_{\mathrm{loc}}^2\delta^S.
\]
The coarse coordinates are orthonormal, overlap bounds \(|e_P|\) by one, and projection contraction follows from the displayed Pythagorean identity. Hence
\[
\|r_u\|_{L_2(\nu_h)}=1,\qquad
\|e_Pr_u\|_{L_2(\nu_h)}\le1,\qquad
\|\Pi_J(e_Pr_u)\|_{L_2(\nu_h)}\le1,\qquad
\|e_Pg_{\mathrm{rem}}\|_{L_2(\nu_h)}
 \le C_{\mathrm{loc}}h^\gamma.
\]
Applying Cauchy--Schwarz to the remaining two integrals therefore gives
\[
|\bar R_u-(Q_{P,h,J}\vartheta)_u|
\le C_{\mathrm{loc}}^2\delta^S+2C_{\mathrm{loc}}h^\gamma
\le
(C_{\mathrm{loc}}^2+2C_{\mathrm{loc}})
(h^\gamma+\delta^S).
\]
Define
\[
C_{\mathrm{bias}}
=\sqrt q\,(C_{\mathrm{loc}}^2+2C_{\mathrm{loc}}).
\]
Squaring the coordinate bounds, summing over the \(q\) coordinates, and taking square roots proves
\[
\|\bar R-Q_{P,h,J}\vartheta\|_2
\le C_{\mathrm{bias}}(h^\gamma+\delta^S).
\]%
% lean: CausalSmith.Stat.TwosamplePointcateAnnotationFrontier.population_bias_coordinate_bound

\item Define the squared deviations, the bias scale, and the good event by
\[
\begin{aligned}
Z_R(\mathcal D_{n,m},U)
 &=\|\widehat R_{h,J}-\bar R\|_2^2,\\
Z_Q(\mathcal D_{n,m},U)
 &=\|\widehat Q_{h,J}-Q_{P,h,J}\|_F^2,\\
b_{\mathrm{bias}}&=h^\gamma+\delta^S,\\
g_\varepsilon&=\frac{\varepsilon(1-\varepsilon)}2,\\
G_{\mathrm{good}}
 &=\{(\mathcal D_{n,m},U):
       \lambda_{\min}(\widehat Q_{h,J})\ge g_\varepsilon\}.
\end{aligned}
\]
The first two functions ignore \(U\).

On \(G_{\mathrm{good}}\), the Rayleigh bound is
\[
\xi_{\mathrm{inv}}^\top\widehat Q_{h,J}\xi_{\mathrm{inv}}
 \ge g_\varepsilon\|\xi_{\mathrm{inv}}\|_2^2
\qquad(\xi_{\mathrm{inv}}\in\mathbb R^q).
\]
It excludes a nonzero kernel and thus makes the matrix invertible.
For arbitrary \(\xi_{\mathrm{rhs}}\in\mathbb R^q\), set
\[
\xi_{\mathrm{inv}}=\widehat Q_{h,J}^{-1}\xi_{\mathrm{rhs}}.
\]
Cauchy--Schwarz then gives
\[
g_\varepsilon\|\xi_{\mathrm{inv}}\|_2^2
\le \xi_{\mathrm{inv}}^\top\xi_{\mathrm{rhs}}
\le \|\xi_{\mathrm{inv}}\|_2\|\xi_{\mathrm{rhs}}\|_2,
\]
and consequently
\[
\|\widehat Q_{h,J}^{-1}\xi_{\mathrm{rhs}}\|_2
\le g_\varepsilon^{-1}\|\xi_{\mathrm{rhs}}\|_2.
\]
For \(\xi_{\mathrm{inv}}=0\) this is immediate; otherwise division by its positive norm proves it.

The inverse identity and the residual decomposition are
\[
\widehat Q_{h,J}^{-1}\widehat R_{h,J}-\vartheta
=
\widehat Q_{h,J}^{-1}
 (\widehat R_{h,J}-\widehat Q_{h,J}\vartheta),
\]
\[
\widehat R_{h,J}-\widehat Q_{h,J}\vartheta
=
(\widehat R_{h,J}-\bar R)
+(\bar R-Q_{P,h,J}\vartheta)
+(Q_{P,h,J}-\widehat Q_{h,J})\vartheta.
\]
The triangle inequality and rowwise Cauchy--Schwarz imply
\[
\|\widehat R_{h,J}-\widehat Q_{h,J}\vartheta\|_2
\le
\sqrt{Z_R}
+C_{\mathrm{bias}}b_{\mathrm{bias}}
+C_{\mathrm{loc}}\sqrt{Z_Q}.
\]
Because both arm means lie in \([0,1]\),
\[
\tau_P(x_0)\in[-1,1].
\]
Clipping toward this interval decreases absolute error: below its lower endpoint the clipped endpoint lies between the value and the target, above its upper endpoint the same holds, and inside the interval clipping preserves the value. Thus, on \(G_{\mathrm{good}}\),
\[
\begin{aligned}
|\widehat T_{h,J}-\tau_P(x_0)|
&\le g_\varepsilon^{-1}\|r_0\|_2
 \|\widehat R_{h,J}-\widehat Q_{h,J}\vartheta\|_2\\
&\le g_\varepsilon^{-1}C_{\mathrm{loc}}
 \left(\sqrt{Z_R}
 +C_{\mathrm{bias}}b_{\mathrm{bias}}
 +C_{\mathrm{loc}}\sqrt{Z_Q}\right).
\end{aligned}
\]
Define the positive finite constant
\[
C_{\mathrm{pt}}
=g_\varepsilon^{-1}C_{\mathrm{loc}}
 (1+C_{\mathrm{bias}}+C_{\mathrm{loc}}).
\]
All three coefficients in the preceding parentheses are at most
\(1+C_{\mathrm{bias}}+C_{\mathrm{loc}}\). On the complementary event the decision is zero and its error is at most \(1\). Combining these branches gives, on every sample,
\[
|\widehat T_{h,J}-\tau_P(x_0)|
\le
C_{\mathrm{pt}}
 (b_{\mathrm{bias}}+\sqrt{Z_R}+\sqrt{Z_Q})
+\mathbf1_{G_{\mathrm{good}}^c}.
\]%
% lean: CausalSmith.Stat.TwosamplePointcateAnnotationFrontier.rectangular_pointwise_error_bound

\item By \cref{obj:lem:population-positivity},
\[
v_{\mathrm{unit}}^\top Q_{P,h,J}v_{\mathrm{unit}}
\ge\varepsilon(1-\varepsilon)=2g_\varepsilon
\qquad
(v_{\mathrm{unit}}\in\mathbb R^q,\ \|v_{\mathrm{unit}}\|_2=1).
\]
For each such vector, finite Cauchy--Schwarz gives
\[
\left|
v_{\mathrm{unit}}^\top
(\widehat Q_{h,J}-Q_{P,h,J})v_{\mathrm{unit}}
\right|^2
\le Z_Q.
\]
If \(Z_Q<g_\varepsilon^2\), every unit-vector quadratic form of
\(\widehat Q_{h,J}\) is at least \(g_\varepsilon\), by subtracting this deviation from \(2g_\varepsilon\). Its infimum is therefore at least \(g_\varepsilon\). Consequently,
\[
G_{\mathrm{good}}^c
\ \Longrightarrow\
Z_Q\ge g_\varepsilon^2,
\qquad
\mathbf1_{G_{\mathrm{good}}^c}
\le g_\varepsilon^{-1}\sqrt{Z_Q}.
\]
Define
\[
\kappa_{\mathrm{mean}}=C_{\mathrm{pt}}+g_\varepsilon^{-1}.
\]
The preceding pointwise bound now becomes
\[
|\widehat T_{h,J}-\tau_P(x_0)|
\le
\kappa_{\mathrm{mean}}
 (b_{\mathrm{bias}}+\sqrt{Z_R}+\sqrt{Z_Q}).
\]

Define the variance scale by
\[
V_{\mathrm{var}}
=\frac1{nh^d}+\frac{qJ^d}{nNh^{2d}}.
\]
By \cref{obj:lem:rectangular-sampling-bound}, there is a uniform positive finite constant \(C_{\mathrm{samp}}\) such that
\[
\mathbb E_{\mathsf E_{n,m}(P)}Z_R
+\mathbb E_{\mathsf E_{n,m}(P)}Z_Q
\le C_{\mathrm{samp}}V_{\mathrm{var}}.
\]
Both deviations are nonnegative. Cauchy--Schwarz under the probability sampling law therefore gives
\[
\begin{aligned}
\mathbb E_{\mathsf E_{n,m}(P)}\sqrt{Z_R}
&\le
\sqrt{\mathbb E_{\mathsf E_{n,m}(P)}Z_R}
\le\sqrt{C_{\mathrm{samp}}}\sqrt{V_{\mathrm{var}}},\\
\mathbb E_{\mathsf E_{n,m}(P)}\sqrt{Z_Q}
&\le
\sqrt{\mathbb E_{\mathsf E_{n,m}(P)}Z_Q}
\le\sqrt{C_{\mathrm{samp}}}\sqrt{V_{\mathrm{var}}}.
\end{aligned}
\]
The decision and target both belong to \([-1,1]\), so their absolute difference is at most \(2\) and is integrable. Integrating the pointwise inequality and defining
\[
C_{\mathrm{mean}}
=\kappa_{\mathrm{mean}}(1+2\sqrt{C_{\mathrm{samp}}})
\]
proves
\[
\begin{aligned}
\mathcal R_{n,m}(\widehat T_{h,J},P)
&\le\kappa_{\mathrm{mean}}
 \left(b_{\mathrm{bias}}
       +2\sqrt{C_{\mathrm{samp}}}\sqrt{V_{\mathrm{var}}}\right)\\
&\le C_{\mathrm{mean}}
 (b_{\mathrm{bias}}+\sqrt{V_{\mathrm{var}}}).
\end{aligned}
\]%
% lean: CausalSmith.Stat.TwosamplePointcateAnnotationFrontier.rectangular_risk_sqrt_variance

\item Define
\[
\rho_{\mathrm{lab}}=(nh^d)^{-1/2},\qquad
\rho_{\mathrm{rect}}=(nNh^d\delta^d)^{-1/2},\qquad
a_{\mathrm{noise}}=\rho_{\mathrm{lab}}+\rho_{\mathrm{rect}},
\qquad
\kappa_{\mathrm{rank}}=1+\sqrt q.
\]
All these quantities are nonnegative, and
\(\kappa_{\mathrm{rank}}\ge1\). Since \(\delta=h/J\),
\[
\rho_{\mathrm{lab}}^2=\frac1{nh^d},\qquad
\rho_{\mathrm{rect}}^2
=\frac{J^d}{nNh^{2d}},\qquad
V_{\mathrm{var}}
=\rho_{\mathrm{lab}}^2+q\rho_{\mathrm{rect}}^2.
\]
It follows that
\[
V_{\mathrm{var}}
\le
(\rho_{\mathrm{lab}}+\sqrt q\,\rho_{\mathrm{rect}})^2,
\qquad
\rho_{\mathrm{lab}}+\sqrt q\,\rho_{\mathrm{rect}}
\le\kappa_{\mathrm{rank}}a_{\mathrm{noise}},
\]
and hence
\[
\sqrt{V_{\mathrm{var}}}
\le\kappa_{\mathrm{rank}}a_{\mathrm{noise}}.
\]
Define
\[
C_{\mathrm{fixed}}
=C_{\mathrm{mean}}\kappa_{\mathrm{rank}}+2.
\]
Using also
\(b_{\mathrm{bias}}\le\kappa_{\mathrm{rank}}b_{\mathrm{bias}}\), the preceding risk bound gives
\[
\mathcal R_{n,m}(\widehat T_{h,J},P)
\le C_{\mathrm{fixed}}
 (b_{\mathrm{bias}}+a_{\mathrm{noise}}).
\]
If \(b_{\mathrm{bias}}+a_{\mathrm{noise}}\le1\), this is the desired bound with the minimum. If
\(b_{\mathrm{bias}}+a_{\mathrm{noise}}>1\), the pointwise diameter bound gives
\[
\mathcal R_{n,m}(\widehat T_{h,J},P)
\le2\le C_{\mathrm{fixed}}.
\]
Thus, uniformly over \(P\in\mathcal M\),
\[
\mathcal R_{n,m}(\widehat T_{h,J},P)
\le C_{\mathrm{fixed}}
\min\left\{1,\,
h^\gamma+\delta^S+(nh^d)^{-1/2}
 +(nNh^d\delta^d)^{-1/2}\right\}.
\]%
% lean: CausalSmith.Stat.TwosamplePointcateAnnotationFrontier.rectangular_fixed_tuning_risk

\item We next establish the public tuning bound. Define
\[
\Delta_{\mathrm{o}}=2\gamma+d,\qquad
\Delta=2\gamma+d+\frac{\gamma d}{S},
\qquad
c_{\mathrm{lab}}=2^{\Delta_{\mathrm{o}}/2},
\qquad
c_{\mathrm{rect}}
=2^{d/2}2^{\Delta/2}+2^{\Delta_{\mathrm{o}}},
\qquad
C_{\mathrm{bal}}=2+c_{\mathrm{lab}}+c_{\mathrm{rect}}.
\]
The parameter domain gives \(S>0\), \(\gamma>0\),
\(\Delta_{\mathrm{o}}>0\), and \(\Delta>0\), so these constants are positive and finite.

For fixed \(n\ge2\) and \(m\ge0\), define the two bandwidth orders
\[
u_{\mathrm{o}}=n^{-1/\Delta_{\mathrm{o}}},
\qquad
u_{\mathrm{rect}}=(nN)^{-1/\Delta}.
\]
Because \(n\ge1\) and \(nN\ge1\),
\[
0<u_{\mathrm{o}}\le1,\qquad
0<u_{\mathrm{rect}}\le1.
\]
The definition in \cref{obj:def:tuning} consequently gives
\[
h_*=\frac12\max\{u_{\mathrm{o}},u_{\mathrm{rect}}\},
\qquad
0<h_*\le\frac12,\qquad
u_{\mathrm{o}}\le2h_*,
\qquad
u_{\mathrm{rect}}\le2h_*.
\]
Moreover,
\[
\begin{aligned}
h_*^\gamma
&\le \max\{u_{\mathrm{o}}^\gamma,u_{\mathrm{rect}}^\gamma\}\\
&=\max\left\{
 n^{-\gamma/\Delta_{\mathrm{o}}},
 (nN)^{-\gamma/\Delta}\right\}
=r_{\mathrm{up}}(n,m),
\end{aligned}
\]
using \cref{obj:def:sharp-rate}.

Taking logarithms of \(u_{\mathrm{o}}\le2h_*\) gives
\[
-\log n
\le\Delta_{\mathrm{o}}(\log2+\log h_*).
\]
Therefore
\[
\begin{aligned}
\log\bigl((nh_*^d)^{-1/2}\bigr)
&=-\frac12(\log n+d\log h_*)\\
&\le\frac{\Delta_{\mathrm{o}}}{2}\log2
   +\gamma\log h_*,
\end{aligned}
\]
where \(\Delta_{\mathrm{o}}-d=2\gamma\). Exponentiating yields
\[
(nh_*^d)^{-1/2}\le c_{\mathrm{lab}}h_*^\gamma.
\]%
% lean: CausalSmith.Stat.TwosamplePointcateAnnotationFrontier.bandwidth_monomial_balance

\item Suppose first that \(S<\gamma\). Define
\[
z_{\mathrm{grid}}=h_*^{-(\gamma/S-1)}.
\]
Then \(z_{\mathrm{grid}}\ge1\), and \cref{obj:def:tuning} gives
\[
J_*=\lceil z_{\mathrm{grid}}\rceil,\qquad
1\le J_*,\qquad
z_{\mathrm{grid}}\le J_*<z_{\mathrm{grid}}+1
\le2z_{\mathrm{grid}}.
\]
Since \(h_*/z_{\mathrm{grid}}=h_*^{\gamma/S}\), these inequalities imply
\[
0<\frac{h_*}{J_*},\qquad
\frac12h_*^{\gamma/S}
\le\frac{h_*}{J_*}\le h_*^{\gamma/S},
\qquad
\left(\frac{h_*}{J_*}\right)^S\le h_*^\gamma.
\]
The lower cell-side bound gives
\[
\log(h_*/J_*)\ge\frac{\gamma}{S}\log h_*-\log2.
\]
Substituting this inequality into the logarithm of the rectangular sampling term and exponentiating proves
\[
\left(nNh_*^d(h_*/J_*)^d\right)^{-1/2}
\le
2^{d/2}
\left(nNh_*^{d+\gamma d/S}\right)^{-1/2}.
\]
Similarly, \(u_{\mathrm{rect}}\le2h_*\) gives
\[
-\log(nN)\le\Delta(\log2+\log h_*).
\]
As \(\Delta-(d+\gamma d/S)=2\gamma\),
\[
\begin{aligned}
\log\left(\left(nNh_*^{d+\gamma d/S}\right)^{-1/2}\right)
&=-\frac12\left(\log(nN)
 +(d+\gamma d/S)\log h_*\right)\\
&\le\frac{\Delta}{2}\log2+\gamma\log h_*.
\end{aligned}
\]
Consequently,
\[
\left(nNh_*^d(h_*/J_*)^d\right)^{-1/2}
\le2^{d/2}2^{\Delta/2}h_*^\gamma
\le c_{\mathrm{rect}}h_*^\gamma.
\]%
% lean: CausalSmith.Stat.TwosamplePointcateAnnotationFrontier.ceil_grid_side_bounds

\item Suppose instead that \(S\ge\gamma\), including equality. The exponent in
\cref{obj:def:tuning} is then zero, so
\[
J_*=1,\qquad
0<\frac{h_*}{J_*}=h_*,
\qquad
\left(\frac{h_*}{J_*}\right)^S
=h_*^S\le h_*^\gamma,
\]
because \(0<h_*\le1\).

The oracle bandwidth inequality also reads
\[
(n^2)^{-1/(2\Delta_{\mathrm{o}})}
=u_{\mathrm{o}}\le2h_*.
\]
Taking logarithms gives
\[
-\log(n^2)
\le2\Delta_{\mathrm{o}}(\log2+\log h_*).
\]
Using \(2\Delta_{\mathrm{o}}-2d=4\gamma\), we obtain
\[
\begin{aligned}
\log\left((n^2h_*^{2d})^{-1/2}\right)
&=-\frac12\left(\log(n^2)+2d\log h_*\right)\\
&\le\Delta_{\mathrm{o}}\log2+2\gamma\log h_*.
\end{aligned}
\]
Since \(N\ge n\) and \(h_*^{2\gamma}\le h_*^\gamma\),
\[
\begin{aligned}
\left(nNh_*^d(h_*/J_*)^d\right)^{-1/2}
&=(nNh_*^{2d})^{-1/2}\\
&\le(n^2h_*^{2d})^{-1/2}\\
&\le2^{\Delta_{\mathrm{o}}}h_*^{2\gamma}\\
&\le c_{\mathrm{rect}}h_*^\gamma.
\end{aligned}
\]
Thus both cases give admissible tuning and the four-term balance
\[
\begin{gathered}
0<h_*\le\frac12,\qquad 1\le J_*,\\
h_*^\gamma+(h_*/J_*)^S+(nh_*^d)^{-1/2}
+\left(nNh_*^d(h_*/J_*)^d\right)^{-1/2}\\
\le(2+c_{\mathrm{lab}}+c_{\mathrm{rect}})h_*^\gamma
\le C_{\mathrm{bal}}r_{\mathrm{up}}(n,m).
\end{gathered}
\]%
% lean: CausalSmith.Stat.TwosamplePointcateAnnotationFrontier.tuning_balance

\item Finally, define one constant for both assertions:
\[
C=C_{\mathrm{fixed}}(C_{\mathrm{bal}}+1).
\]
It is positive and finite, and
\[
C_{\mathrm{fixed}}\le C,\qquad
C_{\mathrm{fixed}}C_{\mathrm{bal}}\le C.
\]
The fixed-tuning envelope has a nonnegative minimum, so enlarging its constant from \(C_{\mathrm{fixed}}\) to \(C\) and taking the supremum over \(P\in\mathcal M\) proves the first risk assertion.

For the second assertion, the tuning bounds make
\((h_*,J_*)\) admissible for the fixed-tuning result. The identity in
\cref{obj:def:tuning} supplies
\[
\widehat T_{\mathrm{up}}=\widehat T_{h_*,J_*},
\]
and hence also its finiteness and Borel measurability. Using the minimum's upper bound by its four-term argument gives, for every \(P\in\mathcal M\),
\[
\begin{aligned}
\mathcal R_{n,m}(\widehat T_{\mathrm{up}},P)
&\le C_{\mathrm{fixed}}
 \left\{h_*^\gamma+(h_*/J_*)^S+(nh_*^d)^{-1/2}
 +\left(nNh_*^d(h_*/J_*)^d\right)^{-1/2}\right\}\\
&\le C_{\mathrm{fixed}}C_{\mathrm{bal}}r_{\mathrm{up}}(n,m)\\
&\le Cr_{\mathrm{up}}(n,m).
\end{aligned}
\]
The rate is nonnegative as the maximum of two positive powers of positive sample counts. All constants are uniform over the primitive class, so taking the supremum proves the tuned assertion.%
% lean: CausalSmith.Stat.TwosamplePointcateAnnotationFrontier.rectangular_upper
\end{enumerate}
\end{proof}
